\section{Métodos implícitos de restricción de política}

\subsection{Filtered/Weighted Imitation Learning}

El baseline más simple: solo imitar las trayectorias o transiciones de \textbf{alta recompensa} en los datos, o ponderar por recompensa.

\subsection{IQL: Implicit Q-Learning}

La idea central de IQL: evitar consultar la función Q en acciones fuera del conjunto de datos.

\textbf{Ajuste de la función V}: usa la pérdida expectile para que $V$ tienda a aprender la parte de comportamiento en los datos con valores Q más altos:
$$\hat{V}(s) \leftarrow \arg\min_{V} \mathbb{E}_{(s,a) \sim \mathcal{D}} \left[\ell_\lambda^2\left(V(s) - \hat{Q}(s, a)\right)\right]$$

donde la pérdida expectile asigna un peso mayor al error positivo ($\lambda > 0.5$), lo que hace que $V$ se incline hacia las mejores acciones de los datos.

\textbf{Ajuste de la función Q}: usa $V$ en lugar de $\max$ para evitar consultas OOD:
$$\hat{Q}(s, a) \leftarrow \arg\min_{Q} \mathbb{E}_{(s,a,s') \sim \mathcal{D}} \left[\left(Q(s,a) - (r + \gamma \hat{V}(s'))\right)^2\right]$$

\subsection{Resumen de la sección}

IQL resuelve el problema de la sobreestimación evitando por completo consultar la función Q en acciones fuera de los datos, y es un método de RL fuera de línea sencillo y elegante.

